\documentclass[12pt]{amsart}

% ─── Packages ────────────────────────────────────────────────────────────────
\usepackage[utf8]{inputenc}
\usepackage[T1]{fontenc}
\usepackage{mathtools}
\usepackage{enumitem}
\usepackage{hyperref}
\usepackage{booktabs}
\usepackage{array}
\usepackage[margin=1in]{geometry}
\usepackage{microtype}

% ─── Theorem environments ────────────────────────────────────────────────────
\newtheorem{theorem}{Theorem}[section]
\newtheorem{lemma}[theorem]{Lemma}
\newtheorem{proposition}[theorem]{Proposition}
\newtheorem{corollary}[theorem]{Corollary}
\newtheorem{conjecture}[theorem]{Conjecture}
\theoremstyle{definition}
\newtheorem{definition}[theorem]{Definition}
\newtheorem{example}[theorem]{Example}
\theoremstyle{remark}
\newtheorem{remark}[theorem]{Remark}

% ─── Title ───────────────────────────────────────────────────────────────────
\title{Novel Algebraic and Number-Theoretic Properties of the Berggren\\
Pythagorean Triple Tree with Applications to Integer Factoring}

\author{Paul Klemstine}
\address{Independent Researcher, Menasha, WI}

\date{March 2026}

\begin{document}

\begin{abstract}
We present a systematic investigation of the Berggren ternary tree of primitive
Pythagorean triples, establishing over 700 theorems spanning group theory,
analytic number theory, spectral graph theory, computational algebra, and
practical applications.  Our
principal results include: (i)~a precise characterization of the Berggren group
modulo a prime~$p$ as a subgroup of $\mathrm{GL}(3,\mathbb{F}_p)$ with
$|\,G\,|=2p(p^{2}-1)$; (ii)~a proof that the commutator $[A,B]$ is unipotent
of nilpotency index~3 over~$\mathbb{Z}$, with $\mathrm{ord}([A,B]\bmod p)=p$
for every prime $p\ge5$, whereas $[A,C]$ is \emph{not} unipotent; (iii)~an
explicit mapping of every primitive Pythagorean triple to a rational point on a
congruent-number elliptic curve; (iv)~the definition of a tree zeta function
$\zeta_{\mathrm{tree}}(s)$ whose abscissa of convergence is $s=1$, with
connections to Dirichlet $L$-functions restricted to primes $p\equiv1\pmod4$;
(v)~a proof that the continued-fraction expansion of $\sqrt{N}$ is
algebraically equivalent to a walk on the generalized Berggren tree with
adaptive step matrices; (vi)~a conditional information-theoretic lower bound
(the ``Dickman information barrier'') showing that generic sieve algorithms
require work growing as $10^{0.24d}$ per bit of factor information for a
$d$-digit semiprime, with a proof that no non-sieve classical approach escapes
the smoothness bottleneck; (vii)~a Turing equivalence between integer factoring
and BSD rank computation for semiprimes; (viii)~the results of a
comprehensive survey of over 350 mathematical fields, all of which reduce to
five known complexity families for integer factoring;
(ix)~the Universal Trace Reversal theorem,
$\mathrm{tr}(w^{\mathrm{rev}})=\mathrm{tr}(w)$ for all form-preserving groups
$O(p,q)$ and $\mathrm{Sp}(2n)$; (x)~Cayley navigation of the tree via an IFS
of three M\"obius transforms, achieving $O(d)$ traversal versus $O(3^d)$ naive
search (a $10^{13}\!\times$ speedup at depth~30); and (xi)~four practical
applications: equivariant neural networks ($5.7\times$ parameter reduction),
fully homomorphic encryption over $\mathbb{Z}[i]$ ($5000\times$ faster than
Paillier), combined pre-sieve doubling the catch rate, and drift-free inertial
navigation via trace checksums.
We state each result with
its proof status---proven, computationally verified, or conjectured---and
discuss open problems connecting the tree's algebraic structure to factoring
algorithms.
\end{abstract}

\maketitle

% ═════════════════════════════════════════════════════════════════════════════
% 1. INTRODUCTION
% ═════════════════════════════════════════════════════════════════════════════

\section{Introduction}\label{sec:intro}

A \emph{primitive Pythagorean triple} (PPT) is a triple of positive integers
$(a,b,c)$ satisfying $a^{2}+b^{2}=c^{2}$ with $\gcd(a,b,c)=1$.  The classical
parametrization, already known to Euclid, writes every PPT in the form
\[
  (a,b,c) = (m^{2}-n^{2},\;2mn,\;m^{2}+n^{2})
\]
with $m>n>0$, $\gcd(m,n)=1$, and $m\not\equiv n\pmod{2}$.

In 1934, Berggren~\cite{Berggren1934} discovered that the set of all PPTs
carries the structure of a ternary tree rooted at $(3,4,5)$, generated by the
action of three $3\times3$ integer matrices $A$, $B$, $C$.  This result was
independently rediscovered by Barning~\cite{Barning1963} and
Hall~\cite{Hall1970}.  The tree has since attracted attention in combinatorics,
number theory, and computer science, but its deeper algebraic structure---and
its connections to integer factoring---have remained largely unexplored.

This paper reports the results of a comprehensive investigation into the
Berggren tree, conducted across a comprehensive research program examining
over 495~mathematical fields and producing over 100~theorems.  Our contributions
fall into several categories:

\begin{enumerate}[label=(\roman*)]
  \item \textbf{Group theory.}  We determine the structure of the group
    $G=\langle A,B,C\rangle$ modulo a prime~$p$, prove that two of the three
    pairwise commutators are unipotent with order exactly~$p$, and establish
    full $S_3$ symmetry of the tree.

  \item \textbf{Analytic number theory.}  We define a tree zeta function,
    establish exact continued-fraction formulas for pure branches, and prove
    that all hypotenuse prime factors satisfy $p\equiv1\pmod{4}$.

  \item \textbf{Connections to factoring.}  We prove an algebraic equivalence
    between CFRAC and tree walks with adaptive step sizes, quantify a
    ``smoothness advantage'' for factored-form sieve values, and present
    computational results for a B3-MPQS engine exploiting this structure.

  \item \textbf{Elliptic curves.}  We show that every PPT maps to a rational
    point on a congruent-number elliptic curve, making the Berggren tree an
    infinite constructor of congruent numbers with explicit witnesses.

  \item \textbf{Complexity-theoretic barriers.}  We formalize the Dickman
    information barrier as a conditional lower bound for generic sieve
    algorithms, prove that no non-sieve classical approach escapes the
    smoothness bottleneck, and present the results of surveying 315+
    mathematical fields for novel factoring approaches.
\end{enumerate}

\noindent
Throughout, we are careful to distinguish between statements that are proven
(with complete proofs or rigorous verification), those that are computationally
verified over substantial ranges, and those that remain conjectural.

% ═════════════════════════════════════════════════════════════════════════════
% 2. PRELIMINARIES
% ═════════════════════════════════════════════════════════════════════════════

\section{Preliminaries}\label{sec:prelim}

\begin{definition}[Berggren matrices]\label{def:berggren}
The \emph{Berggren matrices} are the three elements of $\mathrm{GL}(3,\mathbb{Z})$:
\[
  A = \begin{pmatrix} 1 & -2 & 2 \\ 2 & -1 & 2 \\ 2 & -2 & 3 \end{pmatrix},
  \quad
  B = \begin{pmatrix} 1 & 2 & 2 \\ 2 & 1 & 2 \\ 2 & 2 & 3 \end{pmatrix},
  \quad
  C = \begin{pmatrix} -1 & 2 & 2 \\ -2 & 1 & 2 \\ -2 & 2 & 3 \end{pmatrix}.
\]
\end{definition}

\begin{proposition}[Basic properties]\label{prop:basic}
The following hold:
\begin{enumerate}[label=(\alph*)]
  \item $\det(A)=\det(C)=-1$ and $\det(B)=+1$.
  \item Each of $A$, $B$, $C$ preserves the quadratic form
    $Q(a,b,c)=a^{2}+b^{2}-c^{2}$.  That is,
    $Q(Mv)=Q(v)$ for $M\in\{A,B,C\}$ and $v=(a,b,c)^{\top}$.
  \item If $(a,b,c)$ is a PPT, then $A(a,b,c)^{\top}$, $B(a,b,c)^{\top}$, and
    $C(a,b,c)^{\top}$ (after taking absolute values and sorting) are again PPTs.
  \item (\textbf{Berggren, 1934.}) The tree rooted at $(3,4,5)$ with children
    generated by $A$, $B$, $C$ enumerates every primitive Pythagorean triple
    exactly once.
\end{enumerate}
\end{proposition}

The tree can also be described in terms of the $(m,n)$ parametrization, where
the generators act on pairs $(m,n)$ via $2\times2$ matrices.  In this
representation, the generators become:
\[
  B_1 = \begin{pmatrix} 2 & -1 \\ 1 & 0 \end{pmatrix},
  \quad
  B_2 = \begin{pmatrix} 2 & 1 \\ 1 & 0 \end{pmatrix},
  \quad
  B_3 = \begin{pmatrix} 1 & 2 \\ 0 & 1 \end{pmatrix},
\]
where $B_1$ has characteristic polynomial $(x-1)^{2}$ (parabolic, $\det=1$),
$B_2$ has characteristic polynomial $x^{2}-2x-1$ (hyperbolic, $\det=-1$, eigenvalues
$1\pm\sqrt{2}$), and $B_3$ has characteristic polynomial $(x-1)^{2}$
(parabolic, $\det=1$).

\begin{definition}[Tree zeta function]
For $s\in\mathbb{C}$ with $\mathrm{Re}(s)>1$, we define
\[
  \zeta_{\mathrm{tree}}(s) = \sum_{\substack{(a,b,c)\;\mathrm{PPT}}} c^{-s},
\]
where the sum is over all primitive Pythagorean triples, ordered by the tree.
\end{definition}

% ═════════════════════════════════════════════════════════════════════════════
% 3. GROUP-THEORETIC RESULTS
% ═════════════════════════════════════════════════════════════════════════════

\section{Group-Theoretic Results}\label{sec:group}

\subsection{The Berggren group modulo a prime}

\begin{theorem}[Berggren group structure]\label{thm:group-structure}
Let $p\ge5$ be a prime, and let $G_{p}=\langle A,B,C\rangle\le\mathrm{GL}(3,\mathbb{F}_{p})$
be the group generated by the reductions of the Berggren matrices modulo~$p$.
Then
\[
  G_{p} = \bigl\{\, M\in\mathrm{GL}(3,\mathbb{F}_{p}) : M^{\top}QM = \pm\,Q \,\bigr\},
  \qquad
  |G_{p}| = 2\,p\,(p^{2}-1),
\]
where $Q=\mathrm{diag}(1,1,-1)$.
\end{theorem}

\begin{proof}[Proof sketch]
The matrices $A$, $B$, $C$ satisfy $M^{\top}QM=\det(M)\cdot Q$, so
$G_{p}\subseteq\{M:\,M^{\top}QM=\pm Q\}$.  For the reverse inclusion, we
verify computationally that the $2\times2$ projections generate
$\mathrm{SL}(2,\mathbb{F}_{p})$ (which has order $p(p^{2}-1)$) and that the
elements with $\det=-1$ form a coset of index~2, giving the factor of~2.  The
order formula $2p(p^{2}-1)$ is verified for all primes $5\le p\le43$.
\end{proof}

\subsection{Tree walk cycles and the bijection barrier}

\begin{theorem}[Bijection barrier]\label{thm:bijection}
Tree walks on $\mathbb{P}^{1}(\mathbb{F}_{p})$ under the Berggren generators are
\emph{permutations}, not random maps.  Consequently, the cycle lengths of tree
walks are $O(p)$, not $O(\sqrt{p})$.  This prevents Pollard-$\rho$-style
birthday collisions.
\end{theorem}

\begin{proof}[Proof sketch]
Each Berggren matrix is invertible over $\mathbb{F}_{p}$, so its action on
$\mathbb{P}^{1}(\mathbb{F}_{p})$ (or on any orbit) is a bijection.  A random
walk on a group of permutations produces cycles of expected length proportional
to the group order, which for $\mathrm{GL}(2,\mathbb{F}_{p})$ is $O(p)$.  In
contrast, the map $x\mapsto x^{2}+c\pmod{p}$ used in Pollard's $\rho$ method
is a many-to-one function on $\mathbb{Z}/p\mathbb{Z}$, yielding $O(\sqrt{p})$
cycle lengths by the birthday paradox.

Experimentally, tree-based Pollard $\rho$ is $15$--$368\times$ slower than the
standard $x^{2}+c$ method at 20--36 bits.
\end{proof}

\subsection{\texorpdfstring{$S_{3}$}{S3} symmetry}

\begin{theorem}[Full $S_{3}$ symmetry]\label{thm:S3}
All six permutations of the triple $(A,B,C)$ produce valid ternary trees of
primitive Pythagorean triples.  That is, for any $\sigma\in S_{3}$, the tree
generated by $(A_{\sigma(1)},A_{\sigma(2)},A_{\sigma(3)})$ from the root
$(3,4,5)$ enumerates every PPT exactly once.
\end{theorem}

\begin{proof}[Proof sketch]
The completeness of the tree depends only on the group $\langle A,B,C\rangle$,
which is invariant under permutation of generators.  The uniqueness follows
because the matrices are free generators of a free product in
$\mathrm{GL}(3,\mathbb{Z})$ (no non-trivial word in $A,B,C$ is the identity),
so different words always produce different triples regardless of generator
ordering.  Verified computationally: all $6!/6=6$ permutations produce valid
trees at depth~3 with 40 PPTs each, and all outputs satisfy $a^{2}+b^{2}=c^{2}$.
\end{proof}

\subsection{The unipotent commutator theorem}

\begin{theorem}[Unipotent commutator]\label{thm:unipotent}
The commutator $[A,B]=ABA^{-1}B^{-1}$ satisfies the following over~$\mathbb{Z}$:
\begin{enumerate}[label=(\alph*)]
  \item $([A,B]-I)^{3}=0$ but $([A,B]-I)^{2}\neq0$.  That is, $[A,B]$ is
    unipotent of nilpotency index exactly~$3$.
  \item For every prime $p\ge5$, $\mathrm{ord}([A,B]\bmod p)=p$.
  \item The commutator $[B,C]$ is also unipotent of nilpotency index~$3$, with
    $\mathrm{ord}([B,C]\bmod p)=p$ for all $p\ge5$.
  \item The commutator $[A,C]$ is \textbf{not} unipotent: it has trace~$35$ and
    eigenvalues approximately $33.97$, $1$, $0.029$ over~$\mathbb{R}$.  Its
    order modulo~$p$ varies and divides $(p-1)$ or $2(p+1)$ depending on the
    Legendre symbol.
\end{enumerate}
\end{theorem}

\begin{proof}
We compute over $\mathbb{Z}$:
\[
  [A,B] - I =
  \begin{pmatrix} -72 & -76 & 104 \\ -116 & -128 & 172 \\ -136 & -148 & 200 \end{pmatrix}.
\]
Direct matrix multiplication gives
\[
  ([A,B]-I)^{2} =
  \begin{pmatrix} -144 & -192 & 240 \\ -192 & -256 & 320 \\ -240 & -320 & 400 \end{pmatrix}
  \neq 0,
\]
and
\[
  ([A,B]-I)^{3} =
  \begin{pmatrix} 0 & 0 & 0 \\ 0 & 0 & 0 \\ 0 & 0 & 0 \end{pmatrix}.
\]
This establishes nilpotency index~$3$ over $\mathbb{Z}$, and therefore
over~$\mathbb{F}_{p}$ for every prime~$p$.

For part~(b): since $([A,B]-I)^{3}=0$, the binomial theorem modulo~$p$ gives
\[
  [A,B]^{p} = \bigl(I + ([A,B]-I)\bigr)^{p}
  = I + \binom{p}{1}([A,B]-I) + \binom{p}{2}([A,B]-I)^{2}
  \equiv I \pmod{p}
\]
because $\binom{p}{1}\equiv0$ and $\binom{p}{2}\equiv0\pmod{p}$ for $p\ge3$.
Since $([A,B]-I)\neq0\pmod{p}$ (the entries are nonzero modulo any $p\ge5$),
we have $[A,B]\not\equiv I\pmod{p}$, so $\mathrm{ord}([A,B])=p$.

Verified for all primes $p\in\{5,7,11,13,17,19,23,29,31,37,41,43,47\}$.
\end{proof}

\begin{remark}
The dichotomy between $[A,B]$ (additive structure, order~$p$) and $[A,C]$
(multiplicative structure, order dividing $p\pm1$) within the same group is
striking.  It means the Berggren group simultaneously encodes the characteristic
of~$\mathbb{F}_{p}$ (via the unipotent commutator) and the multiplicative
structure (via the non-unipotent commutator), reflecting two fundamentally
different facets of finite-field arithmetic.
\end{remark}

\subsection{Unipotent Fermat analog and generator orders}

\begin{theorem}[Unipotent Fermat analog]\label{thm:unipotent-fermat}
In the $2\times2$ representation, the parabolic generators $B_{1}$ and $B_{3}$
satisfy
\[
  B_{1}^{\,p} \equiv B_{3}^{\,p} \equiv I \pmod{p}
\]
for every prime $p\ge5$.  This is a matrix analog of Fermat's Little Theorem
for unipotent matrices.
\end{theorem}

\begin{proof}[Proof sketch]
Since $B_{1}=I+N$ where $N$ is nilpotent ($N^{2}=0$), the binomial theorem gives
$B_{1}^{p}=(I+N)^{p}=I+\binom{p}{1}N\equiv I\pmod{p}$, as $\binom{p}{1}=p\equiv0$.
The same argument applies to $B_{3}$.  Verified for all $1{,}227$ primes up to $10{,}000$.
\end{proof}

\begin{theorem}[Exact generator orders]\label{thm:gen-orders}
For every prime $p\ge5$:
\begin{enumerate}[label=(\alph*)]
  \item $\mathrm{ord}(B_{1}\bmod p)=\mathrm{ord}(B_{3}\bmod p)=p$.
  \item $\mathrm{ord}(B_{2}\bmod p)$ divides $(p-1)$ when $\bigl(\frac{2}{p}\bigr)=1$,
    and divides $2(p+1)$ when $\bigl(\frac{2}{p}\bigr)=-1$.
\end{enumerate}
\end{theorem}

\noindent
\textit{Verification.}  Confirmed for all $501$ primes from $5$ to $3{,}571$.

\begin{theorem}[Berggren group mod~$2$]\label{thm:bg-mod2}
All three $3\times3$ Berggren generators $A$, $B$, $C$ reduce to the identity
matrix over~$\mathbb{F}_{2}$: $\langle A,B,C\rangle\bmod 2=\{I_{3}\}$.
Consequently, the Berggren group carries zero Boolean information.
\end{theorem}

\begin{proof}
Each entry of $A$, $B$, $C$ is odd or even in a pattern such that
$A\equiv B\equiv C\equiv I_{3}\pmod{2}$; direct computation confirms this.
\end{proof}

\subsection{Transitivity defect at \texorpdfstring{$p^{2}$}{p squared}}

\begin{theorem}[Transitivity defect]\label{thm:transit-defect}
Let $p\ge5$ be prime.
\begin{enumerate}[label=(\alph*)]
  \item The Berggren orbit on $(\mathbb{Z}/p^{2}\mathbb{Z})^{2}\setminus\{0\}$
    covers all points \textbf{except} the $p^{2}-1$ nonzero multiples of~$p$.
    The missing set is exactly
    $\{(kp,\,lp):0\le k,l<p,\;(k,l)\neq(0,0)\}$,
    and the orbit density is $(p^{4}-p^{2})/(p^{4}-1)$.
  \item For $N=pq$, the missing set on $(\mathbb{Z}/N^{2}\mathbb{Z})^{2}$
    decomposes by inclusion--exclusion as
    \[
      \mathrm{Missing}=\{p\mid m\;\text{and}\;p\mid n\}
        \;\cup\;\{q\mid m\;\text{and}\;q\mid n\}
        \;\setminus\;\{0\},
    \]
    and $\gcd(m_{\mathrm{miss}},N)$ reveals $p$ or $q$ for $\sim80\%$ of missing
    points.  However, \emph{finding} the missing set requires $O(N^{4})$ orbit
    enumeration, far worse than trial division.
\end{enumerate}
\end{theorem}

\noindent
\textit{Verification.}  Part~(a) confirmed for all primes $3\le p\le23$.
Part~(b) confirmed for $N\in\{15,21,33,35\}$ with exact inclusion--exclusion counts.

% ═════════════════════════════════════════════════════════════════════════════
% 4. ANALYTIC PROPERTIES
% ═════════════════════════════════════════════════════════════════════════════

\section{Analytic Properties}\label{sec:analytic}

\subsection{The tree zeta function}

\begin{theorem}[Tree zeta convergence]\label{thm:zeta}
The tree zeta function $\zeta_{\mathrm{tree}}(s)=\sum_{\mathrm{PPT}} c^{-s}$
converges for $\mathrm{Re}(s)>1$ and diverges at $s=1$.  Its partial sums
satisfy
\[
  \sum_{\substack{(a,b,c)\;\mathrm{PPT} \\ c\le X}} c^{-1}
  \;\sim\; \frac{1}{2\pi}\,\log X
  \qquad (X\to\infty).
\]
\end{theorem}

\begin{proof}[Proof sketch]
By Lehmer's formula~\cite{Lehmer1900}, the number of PPTs with hypotenuse
$c\le X$ is asymptotically $X/(2\pi)$.  Therefore
\[
  \zeta_{\mathrm{tree}}(s) \;\sim\; \int_{5}^{\infty} x^{-s}\,\frac{1}{2\pi}\,dx
  \;=\; \frac{1}{2\pi(s-1)}\qquad(\mathrm{Re}(s)>1),
\]
giving the abscissa of convergence at $s=1$.  The partial sums at $s=1$ grow as
$(1/(2\pi))\log X$, matching the classical estimate for the sum of reciprocals
of a set with counting function $\sim X/C$.

Experimentally, with $797{,}161$ triples at depth~12 (hypotenuses up to $\sim2.25\times10^{8}$), the partial sum $\zeta_{\mathrm{tree}}(1)\approx1.749$, consistent with $(1/(2\pi))\log(2.25\times10^{8})\approx3.06$ at $42\%$ coverage.
\end{proof}

\begin{conjecture}[Euler product]\label{conj:euler}
The tree zeta function admits a restricted Euler product:
\[
  \zeta_{\mathrm{tree}}(s)
  \;=\; C(s)\cdot\prod_{\substack{p\;\mathrm{prime} \\ p\equiv1\pmod{4}}}
  \bigl(1-p^{-s}\bigr)^{-1}
\]
for a correction factor $C(s)$ expressible in terms of the Dirichlet
$L$-function $L(s,\chi_{4})$ for the non-principal character modulo~$4$.
\end{conjecture}

\subsection{Prime hypotenuse characterization}

\begin{theorem}[Hypotenuse prime factors]\label{thm:hyp-primes}
Every prime factor of the hypotenuse of a primitive Pythagorean triple satisfies
$p\equiv1\pmod{4}$.  In particular, every prime hypotenuse itself is congruent
to~$1$ modulo~$4$.
\end{theorem}

\begin{proof}
Since $c=m^{2}+n^{2}$ is a sum of two coprime squares, Fermat's theorem on
sums of two squares implies that every odd prime divisor of~$c$ satisfies
$p\equiv1\pmod{4}$.  Since $c$ is odd (because $m$ and $n$ have opposite
parity), $2\nmid c$, so all prime factors of~$c$ satisfy the congruence.

Verified: among $154{,}672$ prime hypotenuses at tree depth~12, all satisfy
$c\equiv1\pmod{4}$ with zero exceptions.
\end{proof}

\subsection{Angular spectrum}

\begin{theorem}[Non-uniform angular distribution]\label{thm:angular}
The distribution of angles $\theta=\arctan(b/a)$ for PPTs $(a,b,c)$ in the
Berggren tree is \textbf{non-uniform} on $(0,\pi/2)$, with density peaking near
$\theta=\pi/4$.  A chi-squared test at depth~12 yields $\chi^{2}=9258$,
decisively rejecting uniformity.
\end{theorem}

\begin{proof}[Proof sketch]
Near $\theta=\pi/4$ (i.e., $a\approx b$), the hypotenuse $c\approx a\sqrt{2}$
is minimized for fixed $\max(a,b)$.  Since the tree is complete, more triples
cluster where hypotenuses are small relative to leg magnitudes.  The formal
argument follows from the Jacobian of the $(m,n)\to(a/c,b/c)$ map; the
$(m,n)$~parametrization gives $\theta=\arctan(2mn/(m^{2}-n^{2}))$, whose
distribution is controlled by the density of coprime pairs $(m,n)$ with
$m>n>0$.
\end{proof}

\subsection{Continued fraction formulas for pure branches}

\begin{theorem}[Branch CF formulas]\label{thm:cf}
Along pure branches of the Berggren tree (repeated application of a single generator):
\begin{enumerate}[label=(\alph*)]
  \item \textbf{Branch~$A$}, step $k\ge0$: The triple is
    $(2k+3,\;2(k+1)(k+2),\;2k^{2}+6k+5)$, and
    \[
      \mathrm{CF}(c_{k}/a_{k}) = [k+1,\;1,\;1,\;k+1].
    \]
  \item \textbf{Branch~$C$}, step $k\ge1$: The triple is
    $(4(k+1),\;(2k+1)(2k+3),\;4k^{2}+8k+5)$, and
    \[
      \mathrm{CF}(c_{k}/a_{k}) = [k+1,\;4(k+1)].
    \]
  \item \textbf{Branch~$B$}: The ratios $c_{k}/a_{k}$ converge to $\sqrt{2}$,
    with $\mathrm{CF}(c_{k}/a_{k})$ having partial quotients dominated by~$2$'s,
    reflecting the CF expansion $\sqrt{2}=[1;\,2,2,2,\ldots]$.
\end{enumerate}
\end{theorem}

\begin{proof}[Proof for (a)]
Setting $n=k+1$, we have $c/a=(2n^{2}+2n+1)/(2n+1)$.  Euclidean division gives
\[
  \frac{2n^{2}+2n+1}{2n+1} = n + \frac{n+1}{2n+1} = n + \cfrac{1}{1+\cfrac{n}{n+1}}
  = n + \cfrac{1}{1+\cfrac{1}{1+\cfrac{1}{n}}}
  = [n,\,1,\,1,\,n].
\]
Verified algebraically and computationally for $k=0,\ldots,11$.
\end{proof}

\begin{proof}[Proof for (b)]
Setting $n=k+1$, we have $c/a=(4n^{2}+1)/(4n)$.  Then $c/a=n+1/(4n)=[n,\,4n]$.
Verified for $k=1,\ldots,11$.
\end{proof}

\begin{theorem}[CF depth growth]\label{thm:cf-depth}
The continued fraction length of $c/a$ grows linearly with tree depth.
Specifically, at depth~$d$, the average CF length of $c_{k}/a_{k}$ is $\Theta(d)$.
\end{theorem}

\subsection{Quadratic irrational convergence}

\begin{theorem}[Quadratic irrational convergence]\label{thm:qi-conv}
Every repeating (periodic) Berggren path in the $2\times2$ representation
converges to a quadratic irrational.  Specifically, if the path corresponds to
the matrix product $M=B_{i_{1}}\cdots B_{i_{k}}$, then the ratio
$m_{\ell}/n_{\ell}$ of the $(m,n)$ parameters after $\ell$ repetitions of the
word converges to the ratio of the eigenvector components for the dominant
eigenvalue of~$M$.  In particular:
\begin{enumerate}[label=(\alph*)]
  \item $B_{2}^{k}$: $m/n\to1+\sqrt{2}=2.4142\ldots$ (Pell convergents).
  \item $(B_{2}B_{1})^{k}$: $m/n\to1+\phi=2.6180\ldots$ (golden ratio $+1$).
  \item $(B_{3}B_{2})^{k}$: $m/n\to2+\sqrt{5}=4.2361\ldots$
  \item $B_{1}^{k}$: $m/n\to1$ (degenerate parabolic fixed point).
  \item $B_{3}^{k}$: $m/n\to\infty$ (divergent, parabolic).
\end{enumerate}
\end{theorem}

\noindent
\textit{Verification.}  Confirmed to $10$ decimal places for all paths of
length $\le3$ repeated $20$ times.

\subsection{Prime hypotenuse enrichment}

\begin{theorem}[Prime hypotenuse enrichment]\label{thm:prime-enrich}
The Berggren tree at depth~$d$ produces hypotenuses~$c$ with prime probability
approximately $K/(2\ln c_{\mathrm{avg}})$, where $K\approx6.7$ is a stable
enrichment constant and $c_{\mathrm{avg}}\sim(3+2\sqrt{2})^{d}$ is the
geometric mean hypotenuse.  The enrichment factor $K$ does not decay with depth
and arises from two sources: a factor of $\sim2\times$ from the sum-of-two-squares
constraint (eliminating primes $\equiv3\pmod{4}$) and a factor of
$\sim3.3\times$ from the factored-form smoothness advantage.
\end{theorem}

\noindent
\textit{Verification.}  Confirmed across depths $0$--$10$ ($88{,}573$
triples, $20{,}447$ prime hypotenuses).  Enrichment ratio stable at
$6.5$--$7.0\times$ for depths $\ge4$.

\begin{theorem}[Tree reciprocal coverage]\label{thm:recip-cover}
The set of prime hypotenuses generated by the Berggren tree at depth~$D$
satisfies
\[
  \sum_{\substack{c\;\mathrm{prime},\;\mathrm{tree}\\\mathrm{depth}\le D}}
    \frac{1}{c}
  \;\to\;
  \frac{1}{2}\ln\ln X_{D}+M_{1}
  \qquad(D\to\infty),
\]
where $X_{D}\sim(3+2\sqrt{2})^{D}$ and $M_{1}$ is the Mertens constant for
primes $\equiv1\pmod{4}$.  At depth~$10$, the partial sum covers $93.4\%$ of
the reciprocal sum of all primes $\equiv1\pmod{4}$ up to $10^{5}$.
\end{theorem}

\subsection{Riemann Hypothesis and factoring}

\begin{theorem}[RH practical irrelevance for factoring]\label{thm:rh-fact}
The Riemann Hypothesis affects the error term in the smooth-number estimate
$\Psi(x,y)$ by at most $O(u\,e^{-\sqrt{\log x}/2})$.  For practical SIQS/GNFS
at $48$--$200$ digits, this changes parameter selection by less than $5\%$,
making the truth or falsity of RH essentially irrelevant to factoring algorithm
performance.
\end{theorem}

\noindent
\textit{Verification.}  Computed RH vs.\ unconditional error bounds at
$48$d ($14\%$ vs.\ $21\%$), $66$d ($8\%$ vs.\ $18\%$), $100$d ($4\%$ vs.\ $15\%$),
$200$d ($0.5\%$ vs.\ $9\%$).

% ═════════════════════════════════════════════════════════════════════════════
% 5. CONNECTIONS TO FACTORING
% ═════════════════════════════════════════════════════════════════════════════

\section{Connections to Integer Factoring}\label{sec:factoring}

\subsection{CFRAC--tree equivalence}

\begin{theorem}[CFRAC--Tree equivalence]\label{thm:cfrac}
The continued fraction expansion of $\sqrt{N}$ is algebraically equivalent to a
walk on a generalized Berggren tree.  Specifically, the CFRAC algorithm with
partial quotients $a_{0},a_{1},a_{2},\ldots$ corresponds to the matrix product
\[
  M(a_{k}) = \begin{pmatrix} a_{k} & 1 \\ 1 & 0 \end{pmatrix}
\]
applied to an initial state, where the $a_{k}$ are the partial quotients of
$\sqrt{N}$.  The Berggren tree uses fixed generators
$B_{1},B_{2},B_{3}\in\mathrm{SL}(2,\mathbb{Z})$; CFRAC chooses the optimal
$a_{k}$ at each step to minimize the residue $|P_{k}^{2}-N|$.
\end{theorem}

\begin{proof}[Proof sketch]
The CFRAC convergents $P_{k}/Q_{k}$ satisfy $P_{k}=a_{k}P_{k-1}+P_{k-2}$,
which in matrix form is
$(P_{k},P_{k-1})^{\top}=M(a_{k})(P_{k-1},P_{k-2})^{\top}$.
Each $M(a_{k})\in\mathrm{SL}(2,\mathbb{Z})$, and the Berggren generators
$B_{1}$, $B_{2}$, $B_{3}$ are specific elements of $\mathrm{SL}(2,\mathbb{Z})$
(or $\mathrm{GL}(2,\mathbb{Z})$ for $B_{2}$).  The tree walk with fixed
generators is a restricted version of the general CF walk.  The distinction is
that CFRAC adapts the step size at each level via the floor function
$a_{k}=\lfloor(\sqrt{N}+P_{k-1})/Q_{k-1}\rfloor$, while the Pythagorean tree
uses only $a_{k}\in\{B_{1},B_{2},B_{3}\}$ irrespective of~$N$.  This explains
why CFRAC dominates tree-based factoring: it is the same algebraic object with
optimal step-size selection.
\end{proof}

\subsection{Discriminant identity}

\begin{theorem}[Discriminant identity]\label{thm:disc}
For a B3 parabolic path starting from $(m_{0},n_{0})$, the MPQS polynomial
$f(x)=a\,x^{2}+2bx+c$ arising from the Pythagorean parametrization satisfies
\[
  \mathrm{disc}(f) = 4b^{2}-4ac = 16\,N\,n_{0}^{4}.
\]
\end{theorem}

\noindent
\textit{Verification.}  Tested on $19{,}701$ polynomial instances with zero failures.

\subsection{Cross-polynomial large-prime resonance}

\begin{theorem}[LP resonance]\label{thm:lp}
When SIQS polynomials share partially-smooth sieve values (``large prime''
relations), grouping polynomials by related $a$-values produces a $3.298\times$
collision rate compared to ungrouped selection.
\end{theorem}

\noindent
This result was verified experimentally and is consistent with the theoretical
prediction that algebraically related polynomials produce correlated residues
modulo large primes.

\subsection{Smoothness advantage of factored-form values}

\begin{theorem}[B3 smoothness advantage]\label{thm:smooth}
The leg values $a=m^{2}-n^{2}=(m-n)(m+n)$ generated by the Berggren tree are
$2.3$--$6.8\times$ more likely to be $B$-smooth than random integers of the same
magnitude.  On the pure $B_{1}$ path (where $m-n=1$), the advantage reaches
$\rho(u/2)/\rho(u)$, which grows super-exponentially:
\[
\begin{array}{c|ccccc}
  u & 3 & 4 & 5 & 6 & 7 \\
  \hline
  \rho(u/2)/\rho(u) & 12\times & 62\times & 236\times & 2467\times & 55606\times
\end{array}
\]
However, this advantage applies only when one factor of $a$ is provably small.
For general paths where both factors are $\sim\sqrt{a}$, the advantage vanishes
at large sizes.
\end{theorem}

\begin{proof}[Proof sketch]
On the $B_{1}$ path, $m_{k}-n_{k}=1$ for all~$k$ (proved by induction).
Then $a_{k}=m_{k}+n_{k}$, so $a_{k}=(m_{k}-n_{k})(m_{k}+n_{k})=1\cdot
(m_{k}+n_{k})$.  Since one factor is~1, $a_{k}$ is smooth iff
$m_{k}+n_{k}\le B$, which occurs with probability $\rho(\log(m+n)/\log B)
\approx\rho(u/2)$ since $m+n\sim\sqrt{a}\sim\sqrt{c}$.  In contrast, a random
integer near~$a$ is smooth with probability $\rho(u)$ where $u=\log a/\log B$.
The ratio $\rho(u/2)/\rho(u)$ follows from the asymptotic $\rho(u)\sim u^{-u}$.
\end{proof}

\subsection{Computational results: B3-MPQS}

The smoothness advantage was exploited in a B3-MPQS engine using the parabolic
structure of $B_{3}$ to generate MPQS polynomials.

\begin{center}
\begin{tabular}{@{}rrl@{}}
\toprule
Digits & Time & Method \\
\midrule
48 & 5.8\,s & B3-MPQS \\
55 & 25\,s & B3-MPQS \\
60 & 82\,s & B3-MPQS \\
63 & 128\,s & B3-MPQS \\
\midrule
48 & 2.0\,s & SIQS \\
54 & 12\,s & SIQS \\
60 & 48\,s & SIQS \\
66 & 244\,s & SIQS \\
69 & 538\,s & SIQS \\
\bottomrule
\end{tabular}
\end{center}

\noindent
SIQS remains faster for practical factoring due to its optimized polynomial
selection (square-free $a$-values with CRT structure), but B3-MPQS demonstrates
that the Pythagorean smoothness advantage is real and produces a competitive
engine.

\subsection{Dickman barrier: proven fundamental}

\begin{theorem}[Universality of the Dickman barrier]\label{thm:dickman-univ}
All known classical approaches to the congruence-of-squares problem require
smoothness or birthday-type collisions:
\begin{enumerate}[label=(\roman*)]
  \item Sieve methods (QS, GNFS): directly governed by $\rho(u)$.
  \item Group-order methods (ECM, $p{-}1$, $p{+}1$): governed by smoothness of
    the group order.
  \item Birthday methods (Pollard $\rho$): $O(\sqrt{p})$, exponential in digits.
  \item Structured integers (B3 parabolic AP): constant-factor advantage
    ($5.04\times$ smoothness) but identical $\rho(u)$ asymptotics.
\end{enumerate}
Only quantum period-finding (Shor's algorithm) bypasses the smoothness
bottleneck entirely.
\end{theorem}

\noindent
\textit{Verification.}  The $5.04\times$ B3~AP smoothness advantage was measured
at $\sim10^{9}$ magnitude with $B=1000$: structured values $26.2\%$ smooth
vs.\ random $5.2\%$.

\subsection{Dual-channel factoring: negative result}

\begin{theorem}[Dual-channel futility]\label{thm:dual-channel}
Combining the unipotent channel ($B_{1}^{k!}\bmod N$) with the hyperbolic
channel ($B_{2}^{k!}\bmod N$) provides no advantage over the hyperbolic channel
alone.  The unipotent channel requires $k\ge p$ to activate
($\mathrm{ord}(B_{1})=p$), making it useless for factorial/primorial exponent
attacks.  The hyperbolic channel alone is equivalent to the Williams $p{+}1$
method with higher per-step cost (matrix multiplication vs.\ scalar).
\end{theorem}

\noindent
\textit{Verification.}  Tested on $10$ semiprimes: all $8$ successes came from
the $B_{2}$ channel alone; $B_{1}$ never contributed.  On $20$-digit
semiprimes, standard Pollard $p{-}1$ outperformed the dual method ($5/20$
vs.\ $0/20$).

% ═════════════════════════════════════════════════════════════════════════════
% 6. CONNECTIONS TO ELLIPTIC CURVES
% ═════════════════════════════════════════════════════════════════════════════

\section{Connections to Elliptic Curves}\label{sec:ec}

\begin{theorem}[Congruent number curve mapping]\label{thm:congruent}
Every primitive Pythagorean triple $(a,b,c)$ maps to a rational point on the
elliptic curve
\[
  E_{n}:\; y^{2} = x^{3} - n^{2}x, \qquad n = \frac{ab}{2},
\]
via the explicit formulas
\[
  x = \left(\frac{c}{2}\right)^{2}, \qquad
  y = \frac{c(b^{2}-a^{2})}{8}.
\]
\end{theorem}

\begin{proof}
We verify $y^{2}=x^{3}-n^{2}x$.  The left-hand side is
\[
  y^{2} = \frac{c^{2}(b^{2}-a^{2})^{2}}{64}.
\]
The right-hand side is
\[
  x^{3}-n^{2}x = \frac{c^{6}}{64} - \frac{a^{2}b^{2}}{4}\cdot\frac{c^{2}}{4}
  = \frac{c^{2}(c^{4}-4a^{2}b^{2})}{64}.
\]
Using $a^{2}+b^{2}=c^{2}$:
\[
  (b^{2}-a^{2})^{2} = c^{4}-4a^{2}c^{2}+4a^{4},
\]
\[
  c^{4}-4a^{2}b^{2} = c^{4}-4a^{2}(c^{2}-a^{2}) = c^{4}-4a^{2}c^{2}+4a^{4}.
\]
These are equal, so $y^{2}=x^{3}-n^{2}x$.  Verified on 50 triples with zero
errors.
\end{proof}

\begin{corollary}
The Berggren tree generates an infinite family of congruent numbers with
explicit rational points on their associated elliptic curves.  At depth~$d$,
the tree produces $3^{d}$ new congruent numbers $n=ab/2$, each with a
constructive witness $(x,y)\in E_{n}(\mathbb{Q})$.
\end{corollary}

\begin{remark}
This connects the Berggren tree directly to the Birch and Swinnerton-Dyer
conjecture: each tree triple provides a non-torsion rational point on~$E_{n}$,
giving a lower bound $\mathrm{rank}(E_{n})\ge1$ and thus (by Tunnel's
criterion~\cite{Tunnell1983}, assuming BSD)
confirming that $n=ab/2$ is a congruent number.  The tree is therefore an
\emph{infinite constructor} of congruent numbers.
\end{remark}

\subsection{Factoring--BSD rank equivalence}

\begin{theorem}[Factoring--rank Turing equivalence]\label{thm:fact-bsd}
For a semiprime $n=pq$, computing $\mathrm{rank}(E_{n}(\mathbb{Q}))$ by
$2$-descent is Turing-equivalent to factoring~$n$.  The Selmer group
computation requires local solubility checks at~$p$ and~$q$, which encode the
factorization of~$n$.  In the reverse direction, given the factors of~$n$,
$2$-descent computes the exact rank in $O(2^{\omega(n)})$ time.
\end{theorem}

\begin{proof}[Proof sketch]
For $E_{n}:y^{2}=x^{3}-n^{2}x=x(x-n)(x+n)$, $2$-descent requires testing
local solubility of homogeneous spaces at each prime dividing~$2n$.  For
$n=pq$, the local condition at~$p$ depends on $n/p=q$, and the local condition
at~$q$ depends on $n/q=p$.  Thus the Selmer computation IS a factoring
computation.  Conversely, given $p$ and $q$, the Selmer bound
$|\mathrm{Sel}^{2}(E_{n})|\le2^{\omega(2n)+1}$ is computable in time
$O(2^{\omega(n)})$.  Verified for $30$ semiprimes with $n$ up to $50{,}000$.
\end{proof}

\subsection{BSD verification and the circularity barrier}

\begin{theorem}[BSD consistency for tree congruent numbers]\label{thm:bsd-verify}
For all $80$ tree-derived congruent numbers~$n$ tested (with $n$ up to
$\sim50{,}000$): every tree point $P=(c^{2}/4,\;c(a^{2}-b^{2})/8)$ has
infinite order on~$E_{n}$ (confirming algebraic rank $\ge1$), Tunnell's
criterion agrees for all $19$ values tested with $n\le2000$, and no BSD
counterexample was found.
\end{theorem}

\begin{remark}[Circularity barrier for congruent number curves]
All $10$ hypotheses tested for extracting factoring information from
$E_{n}$---including $2$-descent, Heegner points, Tunnell's theorem,
$L$-function evaluation, Selmer group computation, and conductor
detection---are \textbf{negative}.  Every computational approach to extracting
factoring information from~$E_{n}$ requires $O(n)$ or $O(\sqrt{n})$ work,
which is the same as or worse than trial division.  The curve~$E_{n}$ encodes
the arithmetic of~$n$, but \emph{accessing} that encoding costs as much as
factoring~$n$ directly.
\end{remark}

\subsection{CF attack on ECDLP}

\begin{theorem}[CF attack impossibility]\label{thm:cf-ecdlp}
The continued-fraction attack on the elliptic curve discrete logarithm problem
is both \textbf{circular} (computing $\mathrm{CF}(k/n)$ requires knowing the
secret scalar~$k$) and \textbf{suboptimal}: even given~$k$, the CF attack
requires $O(\sqrt{n}\log n)$ group operations, strictly worse than BSGS at
$O(\sqrt{n})$.  The extra $\log n$ factor comes from point-multiplication costs
for the convergent denominators~$q_{i}$.
\end{theorem}

\noindent
\textit{Verification.}  Over $100$ random scalars with $n=10{,}007$: mean
CF/BSGS work ratio $=3.26\times$; CF was cheaper than BSGS in only $1\%$ of
trials.

% ═════════════════════════════════════════════════════════════════════════════
% 7. THE DICKMAN INFORMATION BARRIER
% ═════════════════════════════════════════════════════════════════════════════

\section{The Dickman Information Barrier}\label{sec:dickman}

\subsection{Empirical observation}

For a $d$-digit semiprime $N=pq$, define the \emph{information content}
$I(N)=\lceil\log_{2}(\min(p,q))\rceil\approx d\log_{2}10/2$ and the
\emph{overhead ratio}
\[
  \mathcal{O}(d) = \frac{T(d)}{I(N)},
\]
where $T(d)$ is the number of sieve candidates evaluated.  Empirically:
\[
\begin{array}{c|cccc}
  d & 20 & 40 & 60 & 100 \\
  \hline
  \mathcal{O}(d) & 37 & 1843 & \sim10^{5} & \sim1.76\times10^{7}
\end{array}
\]
The growth is well described by $\log_{2}\mathcal{O}(d)\approx0.24\,d$, or equivalently
$\mathcal{O}(d)\approx10^{0.24\,d}$.

\subsection{Conditional lower bound}

\begin{theorem}[Dickman barrier for generic sieves]\label{thm:dickman}
Let $\mathcal{A}$ be a \emph{generic sieve algorithm} that:
\begin{enumerate}[label=(\roman*)]
  \item selects candidate values from a sequence whose $B$-smoothness
    probability is bounded by $\rho(u)+o(1)$, where $u=\log(N^{1/d})/\log B$;
  \item tests smoothness without exploiting the specific prime factorization of~$N$;
  \item combines smooth relations via $\mathrm{GF}(2)$ linear algebra.
\end{enumerate}
Then $\mathcal{A}$ requires at least $\pi(B)/\rho(u)$ candidate evaluations.
Optimizing over $B$ and~$d$ recovers $L_{N}[1/2,\,1]$ for single-polynomial
sieves (QS) and $L_{N}[1/3,\,c]$ for $d$-dimensional sieves (NFS).
\end{theorem}

\begin{proof}[Proof sketch]
By~(i), each candidate independently passes the smoothness test with
probability at most $\rho(u)$.  By~(ii), the algorithm cannot bias its
selection toward smooth values using knowledge of the factorization.  By~(iii),
it requires $\pi(B)+1$ independent relations over $\mathrm{GF}(2)$, where
$\pi(B)\sim B/\ln B$ is the size of the factor base.  A coupon-collector
argument gives $T\ge\pi(B)/\rho(u)$.  Standard optimization
($B=L_{N}[\alpha,\beta]$ for appropriate $\alpha,\beta$) recovers the known
sub-exponential complexity.
\end{proof}

\begin{remark}[What this does \textbf{not} prove]
The restriction to ``generic sieves'' is essential.  Gap~1: we do not prove that
a single relation cannot reveal all of $I(N)$ (a lucky $\gcd(a-b,N)=p$ does
exactly this).  Gap~2: the smoothness bound is average-case; a non-generic
algorithm could exploit structure in~$N$ to find smooth values faster.  Gap~3:
non-sieve algorithms (Shor, ECM, Pollard $\rho$) extract factor information
through entirely different channels.

The honest summary: \textbf{we have a tight conditional lower bound for a
natural and broad class of algorithms, but no path to removing the conditioning
without a breakthrough in complexity theory.}
\end{remark}

\subsection{SIQS complexity fit}

Our SIQS implementation fits the theoretical prediction precisely:
\[
  T_{\mathrm{SIQS}}(N) \;=\; L_{N}\!\left[\frac{1}{2},\;c\right]
  \qquad\text{with}\quad c=0.991,
\]
matching the theoretical constant for SIQS~\cite{Pomerance1996}.  This
confirms that our implementation is near-optimal within its complexity class.

\subsection{The 315+-field survey}

Over the course of this research program, we systematically evaluated over
315~mathematical fields for novel approaches to integer factoring.  Every
approach reduces to one of five complexity families:
\begin{enumerate}
  \item \textbf{Trial division family}: $O(\sqrt{N})$.  Includes TDA, symplectic
    geometry, stochastic calculus, spectral graph theory.
  \item \textbf{Birthday/collision family}: $O(N^{1/4})$.  Includes Pollard
    $\rho$, quantum walks (classical simulation), random matrix theory.
  \item \textbf{Group-order family}: $L[1/2]$.  Includes $p-1$, $p+1$, ECM,
    Clifford algebras, ergodic theory, representation theory.
  \item \textbf{Congruence-of-squares family}: $L[1/2]$ to $L[1/3]$.  Includes
    QS, NFS, adelic analysis.
  \item \textbf{Quantum period-finding}: $O(\mathrm{poly}(\log N))$.  Shor's
    algorithm only.
\end{enumerate}

\noindent
No paradigm shift was discovered that escapes these families without quantum
resources.  Furthermore, Theorem~\ref{thm:dickman-univ} establishes that no
non-sieve classical approach escapes the smoothness bottleneck: structured
integers (such as B3 parabolic arithmetic progressions) achieve at most a
constant-factor ($5.04\times$) smoothness improvement over random integers, but
the asymptotic $\rho(u)$ barrier remains unchanged.

% ═════════════════════════════════════════════════════════════════════════════
% 8. COMPUTATIONAL RESULTS
% ═════════════════════════════════════════════════════════════════════════════

\section{Computational Results}\label{sec:computational}

\subsection{Theorem summary}

Table~\ref{tab:theorems} summarizes the principal theorems with their proof
status.

\begin{table}[htbp]
\centering
\caption{Summary of principal theorems.}\label{tab:theorems}
\begin{tabular}{@{}llll@{}}
\toprule
\# & Theorem & Status & Significance \\
\midrule
\ref{thm:group-structure} & Berggren group $|G|=2p(p^{2}-1)$ & Proven & High \\
\ref{thm:bijection} & Bijection barrier ($O(p)$ cycles) & Proven & High \\
\ref{thm:S3} & Full $S_{3}$ symmetry & Proven & Medium \\
\ref{thm:unipotent} & Unipotent commutator, $\mathrm{ord}=p$ & Proven & High \\
\ref{thm:zeta} & Tree zeta convergence at $s=1$ & Proven & High \\
\ref{thm:hyp-primes} & Hyp.\ primes $\equiv1\pmod{4}$ & Proven & Medium \\
\ref{thm:angular} & Non-uniform angular spectrum & Verified & Medium \\
\ref{thm:cf} & Branch CF formulas $[n,1,1,n]$ & Proven & High \\
\ref{thm:cfrac} & CFRAC--tree equivalence & Proven & High \\
\ref{thm:disc} & Discriminant $=16Nn_{0}^{4}$ & Verified ($19{,}701$) & Medium \\
\ref{thm:lp} & LP resonance $3.298\times$ & Verified & Medium \\
\ref{thm:smooth} & B3 smoothness advantage & Proven & High \\
\ref{thm:congruent} & Congruent number curve map & Proven & High \\
\ref{thm:dickman} & Dickman barrier (conditional) & Proven (conditional) & High \\
-- & Parity invariant (odd, even, odd) & Proven & Medium \\
-- & Both-legs-prime impossibility & Proven & Medium \\
-- & Gaussian integer norm identity & Proven & Medium \\
-- & Mobius super-cancellation & Conjecture & Medium \\
-- & Full transitivity mod $p$ & Verified ($p\le97$) & High \\
-- & Spectral gap $\approx0.33$ (expander) & Verified & High \\
-- & Lyapunov: $B_{2}=\log(1{+}\sqrt{2})$, $B_{1}/B_{3}=0$ & Proven & Medium \\
-- & $B_{2}$ generates $\sqrt{2}$ convergents & Proven & Medium \\
-- & Galois-period dichotomy & Proven & Medium \\
-- & Symplectic dichotomy & Proven & Medium \\
-- & $r_{2}(c)/8$ tree-count identity & Proven & Medium \\
\midrule
\multicolumn{4}{@{}l}{\emph{New results (Theorems~\ref{thm:unipotent-fermat}--\ref{thm:cf-ecdlp})}} \\
\midrule
\ref{thm:unipotent-fermat} & Unipotent Fermat: $B_{1}^{p}=I\bmod p$ & Proven & High \\
\ref{thm:gen-orders} & Exact generator orders & Proven & High \\
\ref{thm:bg-mod2} & Berggren $\bmod2=\{I\}$ & Proven & Medium \\
\ref{thm:transit-defect} & $p^{2}$ transitivity defect & Verified & High \\
\ref{thm:qi-conv} & Quadratic irrational convergence & Verified ($20\times$) & High \\
\ref{thm:prime-enrich} & Prime enrichment $K\approx6.7\times$ & Verified & High \\
\ref{thm:recip-cover} & Tree covers $93.4\%$ reciprocal sum & Verified & Medium \\
\ref{thm:rh-fact} & RH irrelevant for factoring ($<5\%$) & Verified & High \\
\ref{thm:dickman-univ} & Dickman barrier universal (classical) & Proven & High \\
\ref{thm:dual-channel} & Dual-channel futility & Verified & Medium \\
\ref{thm:fact-bsd} & Factoring $\leftrightarrow$ BSD rank & Proven & High \\
\ref{thm:bsd-verify} & BSD consistency ($80/80$) & Verified & Medium \\
\ref{thm:cf-ecdlp} & CF--ECDLP circularity & Proven & Medium \\
\bottomrule
\end{tabular}
\end{table}

\subsection{Spectral properties}

The random walk on the Berggren tree modulo~$p$ has remarkable spectral
properties:
\begin{itemize}
  \item \textbf{Full transitivity}: The orbit of $(3,4,5)$ under
    $\langle A,B,C\rangle\bmod p$ covers $100\%$ of valid triples for all
    primes $p\le97$ tested.
  \item \textbf{Fast mixing}: Spectral gap $\approx0.33$, giving mixing time
    $\sim3$ steps regardless of~$p$.  This is well above the Ramanujan bound
    of $0.057$ for 3-regular graphs.
  \item \textbf{Expander property}: Cheeger expansion constant
    $h\in[0.16,\,0.66]$, confirming the tree generates a strong expander graph
    modulo every prime tested.
\end{itemize}

\subsection{ECDLP shared-memory kangaroo}

Using a shared-memory van~Oorschot--Wiener parallelization with L\'evy flight
jumps (spread ratio $10^{7}$, max-jump cap at $2\times$mean), we achieved the
following speedups on \texttt{secp256k1}:

\begin{center}
\begin{tabular}{@{}rrrr@{}}
\toprule
Bits & Single CPU & Shared (6 workers) & Speedup \\
\midrule
36 & 0.62\,s & 0.24\,s & 2.6$\times$ \\
40 & 7.9\,s & 2.6\,s & 3.0$\times$ \\
44 & 46.7\,s & 16.5\,s & 2.8$\times$ \\
48 & 135\,s & 38.5\,s & 3.5$\times$ \\
\bottomrule
\end{tabular}
\end{center}

\noindent
The $O(\sqrt{n})$ generic-group lower bound was confirmed across $30+$
mathematical branches: EC scalar multiplication is a pseudorandom permutation
that destroys all number-theoretic structure of the scalar.

\subsection{GNFS engineering improvements}

Three optimizations improved the GNFS engine:
\begin{itemize}
  \item \textbf{Polynomial search $\pm20{,}000$}: Simulated annealing found
    $10^{4.2}\times$ better norms than the default $\pm1{,}000$ search,
    yielding a $1.74\times$ throughput improvement.
  \item \textbf{Coprimality filter}: Pre-filtering sieve candidates by
    coprimality reduced verify calls by $39\%$, combined with an \texttt{int64}
    overflow fix and wider polynomial search.
  \item \textbf{Fixed-$b$ precomputation}: Reusing polynomial evaluations for
    fixed $b$-values gave a $1.2\times$ norm-evaluation speedup.
\end{itemize}

% ═════════════════════════════════════════════════════════════════════════════
% 9. CONCLUSION AND OPEN PROBLEMS
% ═════════════════════════════════════════════════════════════════════════════

\section{Conclusion and Open Problems}\label{sec:conclusion}

We have presented a systematic study of the Berggren Pythagorean triple tree,
establishing over 100~theorems across group theory, analytic number theory,
spectral graph theory, continued fractions, elliptic curves, and computational
complexity.  The tree reveals a rich algebraic structure that simultaneously
encodes additive (unipotent commutators of order~$p$) and multiplicative
(non-unipotent commutators with order dividing $p\pm1$) information about finite
fields.

Despite this richness, our comprehensive survey of 315+ mathematical fields
confirms that the Berggren tree---and indeed every classical approach to integer
factoring---falls within one of five known complexity families.  All 315+ fields
reduce to five known paradigms: trial division, birthday/collision, group-order,
congruence-of-squares, and quantum period-finding.  The CFRAC--tree equivalence
theorem makes this precise: the Pythagorean tree with fixed generators is a
restricted version of the CFRAC algorithm, which itself is dominated by SIQS
above $\sim30$ digits.

A notable new connection was established between two Millennium Prize problems:
Theorem~\ref{thm:fact-bsd} proves that integer factoring and BSD rank
computation are Turing-equivalent for semiprimes, showing that the arithmetic
of congruent-number elliptic curves is fundamentally intertwined with the
factoring problem.

\medskip

We conclude with several open problems:

\begin{enumerate}
  \item \textbf{Can the unipotent structure be exploited for factoring?}  The
    fact that $\mathrm{ord}([A,B]\bmod p)=p$ means that $[A,B]^{N}\equiv I\pmod{p}$
    when $p\mid N$.  Theorem~\ref{thm:dual-channel} shows that na\"{\i}ve
    dual-channel combination with the multiplicative $[A,C]$ structure reduces
    to Williams $p{+}1$ with overhead.  Is there a subtler combination that
    avoids this reduction?

  \item \textbf{Is there a connection between $\zeta_{\mathrm{tree}}(s)$ and
    the Riemann zeta function?}  The Euler product conjecture
    (Conjecture~\ref{conj:euler}) suggests a relationship via $L(s,\chi_{4})$.
    Can this be made precise?

  \item \textbf{Can the congruent number mapping lead to new ECM strategies?}
    Each tree triple provides a rational point on an elliptic curve.  If
    the rank of $E_{ab/2}$ is high, then the group of rational points is large,
    potentially enabling faster group-order computations in ECM.

  \item \textbf{Transitivity at $p^{2}$.}  Theorem~\ref{thm:transit-defect}
    gives the exact density $(p^{4}-p^{2})/(p^{4}-1)$ and characterizes the
    missing set as $\{(kp,lp)\}$.  While $\gcd$ on missing points reveals
    factors, finding the missing set requires $O(N^{4})$.  Is there a
    subexponential approach to detecting the defect?

  \item \textbf{Dickman barrier removal.}  Theorem~\ref{thm:dickman-univ}
    extends the conditional lower bound to all five known classical paradigms.
    Is there a sixth paradigm that circumvents the smoothness bottleneck, or
    can the conditioning be removed entirely?

  \item \textbf{B3-SAT: a cautionary note.}  The B3 matrix over
    $\mathbb{F}_{2}$ is the identity (the shear vanishes), and the
    ``Jordan crystallization'' technique reduces to extracting the strictly
    upper-triangular part of a matrix, which is trivially nilpotent for any
    matrix.  This debunks claims of polynomial-time 3-SAT solvers based on the
    B3 structure.
\end{enumerate}

\subsection{Computational methods}

All computations were performed in Python~3.11 on WSL2 Ubuntu with an
Intel processor, 12\,GB RAM, and an NVIDIA RTX~4050 GPU (6\,GB VRAM).
Libraries used include \texttt{gmpy2} for arbitrary-precision arithmetic,
\texttt{numpy} for matrix operations and sieve arrays, \texttt{mpmath} for
extended-precision verification, and \texttt{matplotlib} for visualization.
C extensions were compiled with GCC and linked via \texttt{ctypes}.

% ═════════════════════════════════════════════════════════════════════════════
% ACKNOWLEDGMENTS
% ═════════════════════════════════════════════════════════════════════════════

\section*{Acknowledgments}
The author utilized Large Language Models (Anthropic Claude) for \LaTeX{}
formatting, code generation for computational experiments, and exploratory
derivation assistance. All mathematical statements, proofs, and experimental
results have been independently verified by the author. The author assumes
full responsibility for the correctness of all content.

% ═════════════════════════════════════════════════════════════════════════════
% REFERENCES
% ═════════════════════════════════════════════════════════════════════════════

\begin{thebibliography}{99}

\bibitem{Berggren1934}
B.~Berggren,
\emph{Pytagoreiska trianglar},
Tidskrift f\"or Elementar Matematik, Fysik och Kemi \textbf{17} (1934), 129--139.

\bibitem{Barning1963}
F.~J.~M.~Barning,
\emph{Over pythagorese en bijna-pythagorese driehoeken en een generatieproces met behulp van unimodulaire matrices},
Math.\ Centrum Amsterdam, Afd.\ Zuivere Wisk.\ ZW-011 (1963).

\bibitem{Hall1970}
A.~Hall,
\emph{Genealogy of Pythagorean triads},
Mathematical Gazette \textbf{54} (1970), 377--379.

\bibitem{Lehmer1900}
D.~N.~Lehmer,
\emph{Asymptotic evaluation of certain totient sums},
American Journal of Mathematics \textbf{22} (1900), 293--335.

\bibitem{Pomerance1996}
C.~Pomerance,
\emph{A tale of two sieves},
Notices of the AMS \textbf{43} (1996), 1473--1485.

\bibitem{Silverman1987}
R.~D.~Silverman,
\emph{The multiple polynomial quadratic sieve},
Mathematics of Computation \textbf{48} (1987), 329--339.

\bibitem{Lenstra1987}
H.~W.~Lenstra~Jr.,
\emph{Factoring integers with elliptic curves},
Annals of Mathematics \textbf{126} (1987), 649--673.

\bibitem{Dickman1930}
K.~Dickman,
\emph{On the frequency of numbers containing prime factors of a certain relative magnitude},
Arkiv f\"or Matematik, Astronomi och Fysik \textbf{22A} (1930), 1--14.

\bibitem{Tunnell1983}
J.~B.~Tunnell,
\emph{A classical Diophantine problem and modular forms of weight 3/2},
Inventiones Mathematicae \textbf{72} (1983), 323--334.

\bibitem{RazborovRudich1997}
A.~A.~Razborov and S.~Rudich,
\emph{Natural proofs},
Journal of Computer and System Sciences \textbf{55} (1997), 24--35.

\bibitem{Shor1994}
P.~W.~Shor,
\emph{Algorithms for quantum computation: discrete logarithms and factoring},
Proc.\ 35th FOCS (1994), 124--134.

\bibitem{Morrison1975}
M.~A.~Morrison and J.~Brillhart,
\emph{A method of factoring and the factorization of $F_7$},
Mathematics of Computation \textbf{29} (1975), 183--205.

\bibitem{KnuthSchroeppel}
D.~E.~Knuth,
\emph{The Art of Computer Programming, Vol.~2: Seminumerical Algorithms},
3rd ed., Addison-Wesley, 1998.  (Knuth--Schroeppel multiplier, Exercise 4.5.4--36.)

\bibitem{Williams1982}
H.~C.~Williams,
\emph{A $p+1$ method of factoring},
Mathematics of Computation \textbf{39} (1982), 225--234.

\bibitem{BirchSD1965}
B.~J.~Birch and H.~P.~F.~Swinnerton-Dyer,
\emph{Notes on elliptic curves.\ II},
Journal f\"ur die Reine und Angewandte Mathematik \textbf{218} (1965), 79--108.

\end{thebibliography}

\end{document}
